\section{The exact tensor-residue map to the original spectral trace}
\label{sec:frt}

We retain FTR.1--19 and the coefficient reflection proved in FTR.20--24.
The original tensor pullback is $B(f,g)=R_\chi(f,bg)$, including the
full arithmetic injection and its Taylor units. This section computes
the map from that form to the original reflected multiplication trace.
The construction of the trace of multiplication and its trace pairing
is standard finite-algebra theory: see the Stacks Project Authors,
\emph{The Stacks Project}, Section 49.3, tag
\href{https://stacks.math.columbia.edu/tag/0BVH}{0BVH}, and
Lemma 49.3.1, tag
\href{https://stacks.math.columbia.edu/tag/0BJF}{0BJF}
(version consulted 17 September 2026). That reference also records
the degeneracy caused by nilpotents. The finite calculations below
specialize this established construction to the programme's computed
tensor multiplier; they do not claim to introduce trace forms.

Put $\mathcal E=\mathbb C[S]/\chi$ and retain the full primary
length $e=r+1$, original centers $\lambda_{ab}$, local coordinates
$x=S-\lambda_{ab}$, coefficients $c_{ab,d}$ and indices $d_{ab}$
of FTR.6 and FTR.13. Let $e_{ab}$ denote the complete original
Chinese-remainder idempotents. Define two exact sets of centers by
\[
 \mathcal D=\{(a,b):d_{ab}\geq0\},\qquad
 \mathcal Z=\{(a,b):d_{ab}=-1\},\qquad
 e_{\mathcal D}=\sum_{(a,b)\in\mathcal D}e_{ab},\qquad
 e_{\mathcal Z}=1-e_{\mathcal D}.
 \tag{FRT.1}
\]
Thus $\mathcal E=e_{\mathcal D}\mathcal E\oplus
e_{\mathcal Z}\mathcal E$ with its full original nilpotents.
FTR.20--24 shows that both sets are stable under the original
reflection $(a,b)\mapsto(k-a,b)$. By FTR.16--17,
$\mathcal D$ contains every one of the $4k$ perimeter centers.

\subsection{The calculated multiplier and the exact factorization}

Define the following class in the original algebra, without division
by a possibly singular global $b$:
\[
 a_{\rm tr}=\sum_{(a,b)\in\mathcal D}
       e_{ab}\,\frac{e}{c_{ab,d_{ab}}}
                          (S-\lambda_{ab})^{d_{ab}}.
 \tag{FRT.2}
\]
Every denominator here is nonzero by the definition of $d_{ab}$.
The powers and all coefficients are in the original coordinates.
This explicit multiplier satisfies
\[
 b\,a_{\rm tr}=e_{\mathcal D}\chi'\quad\hbox{in }\mathcal E.
 \tag{FRT.3}
\]
To prove this, fix a center of $\mathcal D$ and put $d=d_{ab}$.
FTR.7 gives
\[
 b(\lambda_{ab}+x)=
 \chi_{ab}(\lambda_{ab}+x)
       \bigl(c_{ab,d}x^{r-d}+
                      \hbox{terms of degree greater than }r-d\bigr)
                         \pmod{x^e}.
\]
Multiplying by $(e/c_{ab,d})x^d$ retains precisely
$e\chi_{ab}(\lambda_{ab})x^r$ modulo $x^e$. Differentiation of
$\chi(\lambda_{ab}+x)=x^e\chi_{ab}(\lambda_{ab}+x)$ gives the
same remainder for $\chi'$. At a center of $\mathcal Z$ both
$b$ and $e_{\mathcal D}\chi'$ are zero. The original
Chinese-remainder isomorphism proves FRT.3 on the whole packet.

Retain the original trace form
\[
 \mathcal W(f,g)=\operatorname{Tr}_{\mathcal E}
                      M_{f^{\sharp_k}g}=R_\chi(f,\chi'g).
 \tag{FRT.4}
\]
For completeness the last identity uses all multiplicities: on a
primary block, multiplication has constant diagonal $f^{\sharp_k}(\lambda)
g(\lambda)$ repeated $e$ times; the local residue of
$\chi'f^{\sharp_k}g/\chi$ has exactly the same value, since
$\chi'/\chi=e/x+$ a function regular at $x=0$.
Thus the trace is the sum of these local residues, as written.
Define its two restrictions, extended by zero on the other summand,
by $\mathcal W_{\mathcal D}(f,g)=\mathcal W(e_{\mathcal D}f,
e_{\mathcal D}g)$ and the corresponding formula for $\mathcal Z$.
The explicit comparison is
\[
 \boxed{\quad
 \mathcal W_{\mathcal D}(f,g)=B(f,a_{\rm tr}g),\qquad
 \mathcal W(f,g)=B(f,a_{\rm tr}g)+\mathcal W_{\mathcal Z}(f,g).
 \quad}\tag{FRT.5}
\]
Indeed $e_{\mathcal D}^{\sharp_k}=e_{\mathcal D}$, and its
idempotence, FRT.3 and FTR.3 give the first equality. Products of
the two complementary reflected idempotents vanish, proving the
second equality with no mixed term removed by assumption.

On any full primary block in $\mathcal Z$, $b=0$ whereas
$\chi'=e\chi_{ab}(\lambda_{ab})x^{e-1}\ne0$. Hence no
multiplication operator on that block can satisfy $b a=\chi'$.
More strongly, $B(f,Tg)$ cannot equal the nonzero restriction of
$\mathcal W$ there for any linear map $T$ on the full packet:
if $f$ is supported at that center, residue duality pairs it only
with the reflected center, which also belongs to $\mathcal Z$;
$bTg$ is zero there. The left side is zero for every $g$, while
the trace form pairs the two nonzero constant value coordinates.
Thus $e_{\mathcal D}\mathcal E$ is the largest direct sum of full
primary factors allowed as the right-input domain when the identity
is tested against every left input $f\in\mathcal E$.
It is also the largest reflection-stable direct sum of full primary
factors on which both restricted forms have this factorization.
Indeed, any additional primary factor in $\mathcal Z$ then brings
its reflected factor, so the same nonzero cross pairing is tested.
A restriction to a non-reflection-stable subspace need not test that
pairing and is not covered by this maximality assertion.
This is the obstruction to that particular
factorization; no assertion about every possible geometric comparison
is inferred from it.

\subsection{All solutions and the induced map on the retained quotient}

On a center in $\mathcal D$, FTR.14 gives
$\operatorname{ann}b=x^{d+1}\mathbb C[x]/x^e$. Consequently
every solution to $b a=\chi'$ on that block is exactly
\[
 a=\frac{e}{c_{ab,d}}x^d+x^{d+1}t,\qquad
 t\in\mathbb C[x]/x^e.
 \tag{FRT.6}
\]
One solution was proved in FRT.3. The difference of two solutions
annihilates $b$, proving that FRT.6 is exhaustive. The class of
this solution modulo the radical of $B$ is therefore independent
of the representative chosen to solve the equation.

Keep that radical as an actual retained submodule, and form the
specified quotient
\[
 \overline{\mathcal E}_B
  =e_{\mathcal D}\mathcal E/\operatorname{rad}(B|_{e_{\mathcal D}\mathcal E}),
 \qquad
 \pi_B:e_{\mathcal D}\mathcal E\longrightarrow\overline{\mathcal E}_B.
 \tag{FRT.7}
\]
This is a named quotient for comparison, not an identification of
its kernel with original theta boundaries. Its local block is
$\mathbb C[x]/x^{d+1}$. The canonical induced map
\[
 \overline T_{\rm tr}:e_{\mathcal D}\mathcal E
          \longrightarrow\overline{\mathcal E}_B,
 \qquad g\longmapsto\pi_B(a_{\rm tr}g)
 \tag{FRT.8}
\]
is, on that same local block, exactly
$g\mapsto(e/c_{ab,d})g(\lambda_{ab})x^d$.
Its image is the direct sum of these original one-dimensional socles;
its kernel is the original nilpotent ideal in
$e_{\mathcal D}\mathcal E$. Both statements follow because
$x^d\ne0$ in $\mathbb C[x]/x^{d+1}$ and its coefficient
$e/c_{ab,d}$ is nonzero. In particular its rank is $|\mathcal D|$,
and the exact sequence is
\[
 0\longrightarrow\sqrt{(0)}_{e_{\mathcal D}\mathcal E}
 \longrightarrow e_{\mathcal D}\mathcal E
 \xrightarrow{\ \overline T_{\rm tr}\ }
 \bigoplus_{(a,b)\in\mathcal D}\mathbb C x^{d_{ab}}
 \longrightarrow0.
 \tag{FRT.9}
\]
The quotient of $B$ by its radical is nondegenerate. FRT.5 becomes
the exact pairing identity
$\mathcal W_{\mathcal D}(f,g)=\overline B(\pi_Bf,
\overline T_{\rm tr}g)$. Its target and kernel are now computed,
rather than assumed to be a faithful or positive realization.

FTR.22 gives
$c_{k-a,b,d}=(-1)^{d+1}\overline{c_{ab,d}}$ and
$d_{k-a,b}=d_{ab}$. Under the original reflection the local
coordinate becomes $-x$ on the reflected block. Substitution in
FRT.2 therefore proves
\[
 a_{\rm tr}^{\sharp_k}=-a_{\rm tr},\qquad
 M_{a_{\rm tr}}A=AM_{a_{\rm tr}}.
 \tag{FRT.10}
\]
For the first identity, the reflected coefficient of
$a_{\rm tr}^{\sharp_k}$ is
$(-1)^d e/\overline{c_{ab,d}}$, which is the negative of
$e/c_{k-a,b,d}$. The second identity is multiplication in the
original commutative algebra. The quotient, socle map and full
factorization are thus compatible with the original action and
its stated reflection, including all nonreduced layers.

\subsection{The actual arithmetic classes under this map}

Let $\mathcal I:\mathcal E\hookrightarrow Q^{\otimes k}$ be the
original full-unit arithmetic map TWB41. Its inverse on its image
is exactly the full Mellin-jet observation followed by division by
$[v_h]_h^{\otimes k}$ and inverse sum substitution. The map
\[
 T_{\rm tr}^{\theta}:
 \mathcal I(e_{\mathcal D}\mathcal E)\longrightarrow
 \mathcal I(e_{\mathcal D}\mathcal E),\qquad
 T_{\rm tr}^{\theta}=\mathcal I M_{a_{\rm tr}}\mathcal I^{-1}
 \tag{FRT.11}
\]
is therefore an explicitly defined arithmetic endomorphism of this
specified image. The original action intertwines with it by FRT.10.
Transporting the computed radical gives exactly the quotient in
FRT.7; no part of that transport asserts that the radical was an
original theta boundary or an absent support.

For the supplied full CRT unit
$s_k=\sum_{a>k/2,b}e_{ab}-\sum_{a<k/2,b}e_{ab}$,
whose full local classes are $+1$ and $-1$, the class
$e_{\mathcal D}s_k$ is a unit in $e_{\mathcal D}\mathcal E$.
Consequently $\overline T_{\rm tr}(e_{\mathcal D}s_k)$ has a
nonzero component in every socle in FRT.9. Its lift
$a_{\rm tr}s_k$ is nonzero at every such full primary block,
and $\mathcal I(a_{\rm tr}s_k)\ne0$ by the actual injection.
The original tensor residue pairs these precise arithmetic classes as
\[
 B(e_{\mathcal D}s_k,a_{\rm tr}s_k)
   =\mathcal W_{\mathcal D}(s_k,s_k)
   =-e|\mathcal D|,\qquad |\mathcal D|\ge4k.
 \tag{FRT.12}
\]
The first equality is the proved comparison; the second uses the
already supplied $s_k^{\sharp_k}s_k=-1$ and the rank
$e|\mathcal D|$ of the specified summand. Thus the earlier negative
trace value now has an explicit preimage pairing in the original
full-unit tensor residue and a calculated action-compatible map.
The pairing is not positive, and the positive canonical boundary
forms retain their separately computed action corrections.
This construction neither assigns a geometric Frobenius to the
packet nor transfers Deligne purity without that comparison.
